% ---------- פרק 5: חלופות תכן, בחירת פתרון וחדשנות ----------
% חישוב המטריצה: א 3.25, ב 2.95, ג 4.25. תרחיש אבטחה 0.30/הסבריות 0.05: א 3.40, ב 2.80, ג 3.95. מפנה: אבטחה ~0.43 (א עוקפת ג), ישימות ~0.45 (ב עוקפת ג).
\chapter{חלופות תכן, בחירת פתרון וחדשנות}\label{ch:design}

הפרק מציג שלוש חלופות שונות בעיקרון, בוחר ביניהן במטריצת החלטה משוקללת, ומפרט את הארכיטקטורה, את מקורות המידע ואת החדשנות של החלופה שנבחרה.

\section{שלוש חלופות שונות בעיקרון}\label{sec:alternatives}

החלופות (טבלה~\ref{tab:alternatives}) נבדלות זו מזו באופן שבו המידע נאסף ומעובד.

{\small
\begin{longtable}{|>{\raggedright\arraybackslash}p{3.6cm}|>{\raggedright\arraybackslash}p{4.8cm}|>{\raggedright\arraybackslash}p{5.1cm}|}
\caption{שלוש חלופות התכן, עם היתרון והמחיר של כל אחת}\label{tab:alternatives}\\
\hline
\textbf{חלופה} & \textbf{היתרון} & \textbf{המחיר} \\ \hline
\endfirsthead
\hline
\textbf{חלופה} & \textbf{היתרון} & \textbf{המחיר} \\ \hline
\endhead
\textbf{א: ניקוד מובנה עם איסוף ידני.} אנליסט ממלא תבנית קבועה, והציון מחושב לפי משקלות \en{AHP} & שקוף, לא תלוי במודל שפה, ולא שולח מידע לשירות חיצוני & איטי ותלוי באנליסט; אין זיהוי פתוח של אירועים ומעקב שיטתי אחר מידע חסר \\ \hline
\textbf{ב: סוכן יחיד עם אחזור מקורות (\en{RAG}).} פרומפט אחד מבקש מהמודל לחפש, לנתח ולהמליץ & מהיר לבנייה וזול להפעלה & קשה לבדוק איך הגיע להמלצה; אין שערים קשיחים, וסיכון ההזיה גבוה \cite{huang2025}; ממד בלי מידע עלול "להיעלם" \\ \hline
\textbf{ג: מערכת מרובת סוכנים עם שערים דטרמיניסטיים.} סוכן לכל ממד, בדיקת סנקציות בקוד ושקלול בנוסחה גלויה & כל ממד מדווח בנפרד על ציון, כיסוי ומקורות; השערים לא תלויים במודל; אפשר להחליף סוכן או מודל בנפרד & מורכבת יותר לבנייה, ויקרה יותר בקריאות למודל \\* \hline
\end{longtable}
}

\section{מטריצת החלטה משוקללת}\label{sec:matrix}

כל קריטריון ומשקלו נגזרו מהדרישות בפרק~\ref{ch:problem} לפני ניקוד החלופות:
\begin{itemize}
\item \textbf{הסבריות ומקור לכל טענה (0.20):} \en{REQ-NFR-01} ו-\en{REQ-FUN-09}.
\item \textbf{טיפול במידע חסר (0.20):} \en{REQ-FUN-03} עד \en{REQ-FUN-05}, לב פער הידע.
\item \textbf{הערכת חוזקות וזיהוי סיכונים (0.20):} \en{REQ-PER-01}, \en{REQ-FUN-06} ו-\en{REQ-FUN-10}.
\item \textbf{אבטחת מידע (0.15):} \en{REQ-CON-01} ו-\en{REQ-CON-03}; כמה מידע על כוונות הרכש יוצא לשירות חיצוני, והאם אפשר להריץ על מודל מקומי.
\item \textbf{ישימות (0.10):} אב-טיפוס עובד בסוף שבוע 10.
\item \textbf{גמישות להתאמה בתעשייה האווירית (0.10):} כמה קל להחליף סוכן, מקור או מודל בלי לשנות את שאר המערכת.
\item \textbf{עלות לדוח (0.05):} משקל נמוך, כי עלות קריאות למודל זניחה ביחס לשעת אנליסט (משוואה~\ref{eq:som}).
\end{itemize}

כל חלופה מקבלת ציון $c_{ij}$ בין 1 ל-5 בכל קריטריון, והציון הכולל של חלופה $j$ הוא:
\begin{equation}\label{eq:pugh}
S_j=\sum_{i=1}^{n} w_i\, c_{ij},\qquad \sum_{i=1}^{n} w_i = 1
\end{equation}

{\small
\begin{longtable}{|>{\raggedright\arraybackslash}p{4.6cm}|c|c|c|c|}
\caption{מטריצת החלטה משוקללת: ציון, ובסוגריים הציון המשוקלל}\label{tab:pugh}\\
\hline
\textbf{קריטריון} & \textbf{משקל} & \textbf{א: ידני} & \textbf{ב: סוכן יחיד} & \textbf{ג: רב-סוכני} \\ \hline
\endfirsthead
\hline
\textbf{קריטריון} & \textbf{משקל} & \textbf{א: ידני} & \textbf{ב: סוכן יחיד} & \textbf{ג: רב-סוכני} \\ \hline
\endhead
הסבריות ומקור לכל טענה & 0.20 & 4 (0.80) & 3 (0.60) & 5 (1.00) \\ \hline
טיפול במידע חסר & 0.20 & 2 (0.40) & 2 (0.40) & 5 (1.00) \\ \hline
הערכת חוזקות וסיכונים & 0.20 & 3 (0.60) & 3 (0.60) & 4 (0.80) \\ \hline
אבטחת מידע & 0.15 & 5 (0.75) & 2 (0.30) & 3 (0.45) \\ \hline
ישימות & 0.10 & 4 (0.40) & 5 (0.50) & 3 (0.30) \\ \hline
גמישות להתאמה & 0.10 & 2 (0.20) & 3 (0.30) & 5 (0.50) \\ \hline
עלות לדוח & 0.05 & 2 (0.10) & 5 (0.25) & 4 (0.20) \\ \hline
\textbf{ציון סופי} & \textbf{1.00} & \textbf{3.25} & \textbf{2.95} & \textbf{4.25} \\* \hline
\end{longtable}
}

\textbf{ביסוס הציונים.}
\begin{itemize}
\item \textbf{הסבריות:} ג' מקבלת 5, כי כל סוכן מחזיר מקור לכל טענה והשערים בקוד; ב' 3, כי מודל יחיד עלול להפיק טענות בלי מקור \cite{huang2025}; א' 4, כי אנליסט מתעד מקורות, אבל לא באופן אחיד.
\item \textbf{מידע חסר:} ג' מקבלת 5, כי כל סוכן חייב להחזיר ציון כיסוי; בגישת הסוכן היחיד הניתוח הוא ברמת הכתבה, בלי מדד לחוסר \cite{zhao2025}; גם בחלופה א' אין מנגנון שיטתי.
\item \textbf{חוזקות וסיכונים:} ג' מקבלת 4, כי יש סוכן לכל ממד, אבל איכות נמדדת באותות עקיפים בלבד.
\item \textbf{אבטחת מידע:} א' מקבלת 5, כי אינה שולחת דבר למודל שפה; ב' 2, כי כל הבקשה, ההקשר והמסקנות עוברים במודל חיצוני אחד; ג' 3, כי היא שולחת למודל רק את שם הספק והרכיב, השער מחושב בקוד בלי מודל, וכל סוכן יכול לעבור בנפרד למודל מקומי.
\item \textbf{ישימות:} ב' מקבלת 5 (פרומפט אחד); א' 4 (בלי פיתוח, אבל מילוי ידני לכל מקרה); ג' 3 (שבעה סוכנים, שער ושקלול).
\item \textbf{גמישות:} ג' מקבלת 5, כי אפשר להחליף סוכן או מודל בלי לגעת בשאר; ב' 3, כי כל שינוי הוא שינוי בפרומפט ארוך אחד; א' 2, כי כל שינוי דורש הדרכה מחדש.
\item \textbf{עלות:} א' מקבלת 2 (שעות אנליסט); ב' 5 (קריאה אחת); ג' 4 (שבע קריאות לפחות, אך לפי \en{AlMahri} ועמיתיו כ-8 סנטים לניתוח \cite{almahri2026}).
\end{itemize}

\textbf{ניתוח רגישות.} בתרחיש שבו אבטחת המידע חשובה הרבה יותר, משקלה עולה ל-0.30 ומשקל ההסבריות יורד ל-0.05; גם אז ג' מובילה, עם 3.95 מול 3.40 לחלופה א' ו-2.80 לחלופה ב'. כשמעלים משקל של קריטריון אחד ומקטינים את השאר באופן יחסי, חלופה א' עוקפת את ג' רק כשמשקל אבטחת המידע עולה על כ-0.43, וחלופה ב' עוקפת את ג' רק כשמשקל הישימות עולה על כ-0.45. גם אם "טיפול במידע חסר" מאבד את כל משקלו, ג' מובילה (4.06). \textbf{ההחלטה: חלופה ג', מערכת מרובת סוכנים עם שערים דטרמיניסטיים.}

\section{ארכיטקטורת המערכת}\label{sec:architecture}

איור~\ref{fig:arch} מציג את הזרימה מבקשת הבדיקה ועד ההמלצה. העיקרון המרכזי הוא \textbf{הפרדה בין מה שמחושב בקוד לבין מה שמודל השפה מעריך}: בדיקת הסנקציות, השקלול וכללי ההחלטה מחושבים בקוד, ולכן ניתנים לבדיקה וחוזרים על עצמם; מודל השפה משמש רק לקריאת טקסט ולהערכה בתוך כל ממד; וההחלטה הסופית נשארת אצל אדם (\en{REQ-NFR-02}).

\begin{figure}[htbp]
\centering
\begin{tikzpicture}[
  box/.style={draw, rounded corners, align=center, font=\footnotesize, minimum height=9mm, inner sep=3pt},
  code/.style={box, fill=gray!15},
  llm/.style={box, fill=blue!8},
  sgA/.style={box, fill=teal!15, text width=1.42cm},
  rkA/.style={box, fill=orange!20, text width=1.42cm},
  human/.style={box, fill=green!12},
  arr/.style={-{Stealth[length=2mm]}, thick},
  node distance=4mm and 5mm]
\node[human, text width=2.8cm] (in) {\foreignlanguage{hebrew}{בקשה: שם הספק, מדינה ורכיב}};
\node[code, text width=2.8cm, below=of in] (er) {\foreignlanguage{hebrew}{זיהוי ישות: שם רשמי, בעלים ורמת מידע \en{A/B/C}}};
\node[code, text width=2.8cm, below=of er] (gate) {\foreignlanguage{hebrew}{שער: סנקציות, בקרת יצוא ומדינה עוינת}};
\node[box, fill=red!12, text width=2.2cm, left=12mm of gate] (rej) {\foreignlanguage{hebrew}{פסילה בשער}};
\node[rkA, below=10mm of gate, xshift=-57mm] (a7) {\foreignlanguage{hebrew}{סוכן סייבר (פסיבי)}};
\node[rkA, right=2mm of a7] (a6) {\foreignlanguage{hebrew}{סוכן אירועים (פתוח)}};
\node[rkA, right=2mm of a6] (a5) {\foreignlanguage{hebrew}{סוכן גיאו\-פוליטי}};
\node[sgA, right=2mm of a5] (a4) {\foreignlanguage{hebrew}{סוכן פיננסי}};
\node[sgA, right=2mm of a4] (a3) {\foreignlanguage{hebrew}{סוכן עמידה ביעדים}};
\node[sgA, right=2mm of a3] (a2) {\foreignlanguage{hebrew}{סוכן התאמה טכנית}};
\node[sgA, right=2mm of a2] (a1) {\foreignlanguage{hebrew}{סוכן איכות ואמינות}};
\node[code, text width=7.4cm, below=10mm of a4] (agg) {\foreignlanguage{hebrew}{שקלול: פרופיל חוזקות $P$, דגלי סיכון $R$, כיסוי $C$, כללי חומרה ושער סייבר מותנה}};
\node[llm, text width=7.4cm, below=of agg] (rep) {\foreignlanguage{hebrew}{כותב הדוח: חוזקות וסיכונים לכל ממד, מקור לכל טענה ורשימת בדיקה ידנית}};
\node[human, text width=7.4cm, below=of rep] (hum) {\foreignlanguage{hebrew}{בודק אנושי: ההחלטה הסופית}};
\draw[arr] (in) -- (er);
\draw[arr] (er) -- (gate);
\draw[arr] (gate) -- (rej);
\foreach \n in {a1,a2,a3,a4,a5,a6,a7}{ \draw[arr] (gate.south) -- (\n.north); }
\foreach \n in {a1,a2,a3,a4,a5,a6,a7}{ \draw[arr] (\n.south) -- (agg.north); }
\draw[arr] (agg) -- (rep);
\draw[arr] (rep) -- (hum);
\end{tikzpicture}

\vspace{1mm}
{\footnotesize \colorbox{teal!15}{סוכן חוזקות} \ \colorbox{orange!20}{סוכן סיכון} \ \colorbox{gray!15}{קוד} \ \colorbox{blue!8}{מודל שפה} \ \colorbox{green!12}{אדם}}
\caption{ארכיטקטורת מצפן ספקים: שער חסימה, שבעה סוכנים במקביל (ארבעה לחוזקות ושלושה לסיכונים), שקלול בקוד והחלטה אנושית}\label{fig:arch}
\end{figure}

\textbf{הפלט של כל סוכן.} נסמן ב-$d$ ממד כלשהו מבין השבעה. הסוכן של ממד $d$ מחזיר ציון $s_d$ בין 0 ל-100 (בשני החלקים ציון גבוה הוא טוב: ספק חזק, או סיכון נמוך), ציון כיסוי $c_d$ בין 0 ל-1 (החלק המשוקלל מהקטגוריות בממד שנבדק בפועל), רשימת טענות עם מקור לכל אחת, ורשימת נושאים שלא נמצא עליהם מידע. טענה בלי מקור אינה נכנסת לציון. סוכן האירועים מוסיף לכל אירוע רמת חומרה (1 עד 4) ורמת ביטחון, ואירוע שאינו מתאים לאף קטגוריה מסומן "אחר / חדש" ומחכה לאישור אדם.

\textbf{השקלול.} $D^{+}$ הם ממדי \textbf{פרופיל החוזקות} (איכות ואמינות, התאמה טכנית, עמידה ביעדים, מצב פיננסי), ו-$D^{-}$ ממדי \textbf{דגלי הסיכון} (מצב גיאופוליטי, אירועים חריגים, אבטחת סייבר). $D^{*}$ היא קבוצת הממדים שבהם $c_d \ge c_{\min}$. המשקלות $w_d$ נקבעים ב-\en{AHP} \textbf{בתוך כל חלק}, וסכומם 1 בכל חלק:
\begin{equation}\label{eq:score}
P=\frac{\sum_{d\in D^{+}\cap D^{*}} w_d\, s_d}{\sum_{d\in D^{+}\cap D^{*}} w_d},\qquad
R=\frac{\sum_{d\in D^{-}\cap D^{*}} w_d\, s_d}{\sum_{d\in D^{-}\cap D^{*}} w_d},\qquad
C=\tfrac{1}{2}\Bigl(\sum_{d\in D^{+}} w_d\, c_d+\sum_{d\in D^{-}} w_d\, c_d\Bigr)
\end{equation}
$R$ הוא ציון דגלי הסיכון, וציון גבוה בו פירושו מעט דגלים. שני חלקים נפרדים מונעים מאיכות גבוהה "לפצות" על סיכון גבוה, ומקטינים את ה-\en{AHP} מ-21 השוואות בזוגות (שבעה ממדים יחד) ל-9 (מטריצה של ארבעה ממדים ומטריצה של שלושה). המכנים מחלקים מחדש, באופן גלוי בדוח, את משקל הממדים שאין עליהם מספיק מידע, ו-$C$ מונע מהיעדר מידע להתפרש כהיעדר סיכון או כאיכות טובה.

\textbf{כללי ההחלטה,} לפי הסדר:
\begin{enumerate}
\item ספק ברשימת סנקציות או בקרת יצוא, או שהוא או בעלי השליטה בו ממדינה עוינת: \textbf{לא מומלץ (פסילה)}, לפני כל חישוב.
\item $C<0.4$, או חלק שאין בו אף ממד עם מספיק מידע: \textbf{מידע לא מספיק}, עם רשימת בדיקה ידנית ושאלון לספק. אם בכל זאת נמצא אירוע ברמה 2 שלא נסגר, ההמלצה היא "\textbf{לא מומלץ}" לפי כלל 4, והאירוע מוצג בראש הדוח.
\item לפי החלק החלש, $m=\min(P,R)$: אם $m\ge 70$ \textbf{מומלץ}; אם $50\le m<70$ \textbf{מומלץ בתנאים}; אם $m<50$ \textbf{לא מומלץ}.
\item \textbf{תקרה:} אירוע ברמת חומרה 2 (כשל אספקה) שלא נסגר מוריד את ההמלצה ל"\textbf{לא מומלץ}" עד לבדיקה ידנית. אירוע סייבר חמור שלא נסגר ולא עצר ייצור מוריד "מומלץ" ל"\textbf{מומלץ בתנאים}" עד לבדיקת אבטחה ידנית. המלצה נמוכה יותר נשארת כמו שהיא.
\end{enumerate}
\textbf{רמת הביטחון} גבוהה אם $C\ge 0.75$ ובינונית אם $0.4\le C<0.75$. הערכים $c_{\min}$, 0.4, 50, 70 ו-0.75 הם ערכי פתיחה: הם נקבעים סופית על \textbf{קבוצת כיול} מתוך מאגר התיקוף, ו-\en{H1} ו-\en{H2} נבדקות על \textbf{קבוצת בדיקה} נפרדת. אם המאגר קטן מכדי לחלק אותו, הספים נשארים ערכי הפתיחה ומדווח ניתוח רגישות שלהם.

\section{הדאטה: מקורות מידע ציבוריים}\label{sec:data}

כל המקורות ציבוריים (\en{REQ-CON-01}), והאיסוף פסיבי בלבד, בלי סריקת מערכות הספק (\en{\mbox{REQ-CON-02}}). טבלה~\ref{tab:data} מפרטת את המקורות לכל ממד. לפני השימוש בכל מקור נבדקים תנאי השימוש, אופן הגישה (ממשק תכנות, \en{API}, או הורדה) והכיסוי הגיאוגרפי (משימה \en{T1.2}, פרק~\ref{ch:workplan}).

{\small
\begin{longtable}{|>{\raggedright\arraybackslash}p{2.2cm}|>{\raggedright\arraybackslash}p{5.8cm}|>{\raggedright\arraybackslash}p{3.8cm}|>{\raggedright\arraybackslash}p{1.6cm}|}
\caption{מקורות מידע ציבוריים מועמדים לכל ממד}\label{tab:data}\\
\hline
\textbf{ממד} & \textbf{מקור מועמד} & \textbf{מה הוא נותן} & \textbf{סוג החברה} \\ \hline
\endfirsthead
\hline
\textbf{ממד} & \textbf{מקור מועמד} & \textbf{מה הוא נותן} & \textbf{סוג החברה} \\ \hline
\endhead
שער חסימה & \en{SDN} של \en{OFAC} \cite{ofacsdn}, \en{Entity List} של \en{BIS} \cite{bis2025}, רשימות האיחוד האירופי, \en{OpenSanctions}, ורשימת המדינות העוינות לפי ההגדרה בפרק~\ref{ch:problem}, מול מקור רשמי שייקבע במשימה \en{T2.1} & האם הישות או בעליה ברשימה, והאם מדינת המושב עוינת או בלי יחסים עם ישראל & כל ספק \\ \hline
גיאופוליטי & מדד \en{GPR} \cite{caldara2022}, מדדי הממשל של הבנק העולמי (\en{WGI}), מרשמי חברות & סיכון מדינה, יחסים עם ישראל ובעלות & ציבורית ופרטית \\ \hline
פיננסי & \en{SEC EDGAR} (\en{10-K}, \en{10-Q}), דוחות בבורסות אחרות, מרשמי חברות שמחייבים דיווח & יחסים פיננסיים ומדד \en{Z} \cite{altman1968}; בפרטית: מדדים עקיפים & ציבורית; פרטית חלקית \\ \hline
עמידה ביעדים & תחזיות בהודעות ובדוחות מול הביצוע & הפער בין ההבטחה לתוצאה & בעיקר ציבורית \\ \hline
אירועים חריגים & חדשות (למשל \en{GDELT}), הודעות רגולטורים, רשומות בתי משפט & אירועים שליליים עם תאריך ומקור & ציבורית ופרטית \\ \hline
סייבר & דיווחי \en{8-K} סעיף 1.05 \cite{sec2023}, הודעות \en{CERT} ורגולטורים, מאגרי \en{ISO 27001} & אירועים שדווחו והסמכות & ציבורית ופרטית \\ \hline
איכות ואמינות & \en{OASIS} של \en{IAQG} (\en{AS9100}), \en{USAspending.gov}, פרסי ספקים, הוראות תקינות (\en{Airworthiness Directives}) & הסמכות, לקוחות, הכרה פומבית, ריקולים & ציבורית ופרטית \\ \hline
התאמה טכנית & מפרטים באתר החברה, פטנטים, הסמכות תהליך (\en{Nadcap}) & התאמת חומרים, תהליכים ויכולות לרכיב & ציבורית ופרטית \\* \hline
\end{longtable}
}

\textbf{מאגר התיקוף.} המתודולוגיה תיבדק לאחור על מקרים ציבוריים לפי סולם החומרה (פרק~\ref{ch:problem}): ברמה 2, \en{Incora} (פשיטת רגל, 2023) \cite{incora2023} ו-\en{ASCO} (כופרה, 2019) \cite{certeu2019}; ברמה 1, חברות שנוספו לרשימות סנקציות ובקרת יצוא בגלל העברת רכיבי תעופה לאיראן \cite{treasury2023,bis2025}; וברמה 3, \en{FACC}, יצרנית אוסטרית של חלקי מטוסים שאיבדה 42 מיליון אירו ב-2016 בהונאת "מנכ"ל מזויף" \cite{reuters2016facc}, כדי לבדוק שאות סיכון אינו מסווג ככשל. מולם תיבחר קבוצה של 5 עד 10 ספקים מוכחים, שמהם מצופה "מומלץ". בניית המאגר המלא היא משימה \en{T1.1}.

לכל מקרה נכתבים מראש התוצאה הצפויה ותאריכי החיתוך. \textbf{כשל שאפשר לצפות}, כמו \en{Incora}, שקדמו לו ירידה בהכנסות ונטל חוב, משמש לבדיקת החיזוי ולמדידת \textbf{זמן ההתרעה}: מספר החודשים בין תאריך החיתוך המוקדם ביותר שבו ההמלצה היא "לא מומלץ" לבין מועד הכשל (שאלת משנה 4). זה אותו רף שלפיו נקבע הסיווג הנכון בפרק~\ref{ch:problem}; התאריך שבו ההמלצה ירדה לראשונה מ"מומלץ" מדווח בנפרד, כאות מוקדם בלבד. \textbf{כשל שאי אפשר לצפות}, כמו מתקפת הכופרה על \en{ASCO}, משמש רק לבדיקת סיווג החומרה. טבלה~\ref{tab:cases} מפרטת לכל סוג מקרה מתי הוא מורץ, מה התוצאה הצפויה, והאם הוא נכנס לחישוב הסיווג הנכון של \en{H1}.

{\small
\begin{longtable}{|>{\raggedright\arraybackslash}p{3.3cm}|>{\raggedright\arraybackslash}p{3.6cm}|>{\raggedright\arraybackslash}p{3.9cm}|>{\raggedright\arraybackslash}p{3.3cm}|}
\caption{התוצאה הצפויה לכל סוג מקרה במאגר התיקוף}\label{tab:cases}\\
\hline
\textbf{המקרה} & \textbf{תאריך החיתוך} & \textbf{התוצאה הצפויה} & \textbf{חישוב \en{H1}} \\ \hline
\endfirsthead
\hline
\textbf{המקרה} & \textbf{תאריך החיתוך} & \textbf{התוצאה הצפויה} & \textbf{חישוב \en{H1}} \\ \hline
\endhead
\en{Incora}: כשל שאפשר לצפות & 12 חודשים לפני פשיטת הרגל, ועוד שני תאריכים מוקדמים יותר & "לא מומלץ" & כן; נמדד גם זמן ההתרעה \\ \hline
\en{ASCO}: כשל שאי אפשר לצפות & אחרי פרסום המתקפה & אירוע ברמת חומרה 2, ולפי כלל 4 "לא מומלץ" & לא; בודק את סיווג החומרה (\en{\mbox{REQ-FUN-07}}) \\ \hline
ישויות ברשימות סנקציות ובקרת יצוא & אחרי הכניסה לרשימה & "פסילה" בשער & לא; בודק את שער החסימה (\en{\mbox{REQ-FUN-01}}) \\ \hline
\en{FACC}: אות סיכון & אחרי פרסום ההונאה & אירוע ברמת חומרה 3; ההמלצה אינה יורדת בגללו ל"לא מומלץ" & לא; בודק שאות סיכון אינו מסווג ככשל \\ \hline
ספקים מוכחים & מועד ההרצה & "מומלץ" & כן \\* \hline
\end{longtable}
}

\textbf{מניעת "ידיעת העתיד".} מודל שפה כבר יודע ש-\en{Incora} פשטה רגל. לכן כל הרצה נעשית \textbf{לפי נקודת זמן}: הסוכנים מקבלים רק מסמכים מלפני תאריך החיתוך, ההנחיה אוסרת ידע מחוץ למסמכים, וכל מקור שצוטט נבדק מול תאריך החיתוך. המקרים שלעיל קדמו לתאריך הידע של המודלים, ולכן זו ההגנה היחידה עליהם; מקרים מ-2025 ואילך יתווספו אם יימצאו (סיכון \en{\mbox{VAL-02}}, פרק~\ref{ch:risks}).

\section{החדשנות והערך המוסף}\label{sec:innovation}

החדשנות אינה במערכת ניטור חדשה, אלא בשלוש נקודות שלא נמצאו בספרות (פרק~\ref{ch:literature}) ובתיעוד המוצרים (פרק~\ref{ch:market}):
\begin{enumerate}
\item \textbf{הערכה בשני חלקים עם ציון כיסוי:} ההמלצה עונה על "כמה הספק טוב", "כמה הוא מסוכן" ו"כמה ידוע עליו", כך שספק פרטי לא נראה בטוח רק כי לא נמצאו עליו חדשות רעות.
\item \textbf{חומרה לפי רציפות האספקה, עם קטגוריה פתוחה:} אירוע מדורג לפי השפעתו על האספקה, ואירוע מסוג חדש אינו נכפה על טקסונומיה קיימת.
\item \textbf{שערים בקוד, הערכה במודל, החלטה אצל אדם:} כל המלצה ניתנת לבדיקה, ואפשר להעביר את המערכת למודל מקומי ברשת הארגון בלי לשנות את הלוגיקה (\en{REQ-CON-03}).
\end{enumerate}
